\documentclass[11pt]{article}
\usepackage[a4paper,margin=1in]{geometry}
\usepackage[T1]{fontenc}
\usepackage[utf8]{inputenc}
\usepackage{lmodern}
\usepackage{enumitem}
\usepackage{listings}
\usepackage{xcolor}
\lstset{basicstyle=\ttfamily\small,breaklines=true,frame=single,backgroundcolor=\color{gray!5}}
\title{E1402 Final Assignment - Full Detailed Answer}
\date{}
\begin{document}
\maketitle

\section*{PART I - Multiple Choice \& Theory (Full + Answers)}
\begin{enumerate}[leftmargin=*]
\item OpenCV default color format: \textbf{b) BGR}.\\
Giai thich: OpenCV doc anh theo thu tu Blue-Green-Red.
\item If \texttt{cv2.imread()} cannot find file: \textbf{c) None}.
\item BGR to grayscale constant: \textbf{c) \texttt{cv2.COLOR\_BGR2GRAY}}.
\item Filter for Gaussian noise: \textbf{d) Gaussian Blur}.
\item MaxPooling(2x2): \textbf{b)} reduce spatial resolution by keeping maxima.
\item Why \texttt{waitKey()} after \texttt{imshow()}: \textbf{b)} render window and wait for key press.
\end{enumerate}

\subsection*{7) Contour Detection}
\begin{itemize}[leftmargin=*]
\item \texttt{cv2.findContours()}: detect contours in binary image.
\item \texttt{cv2.drawContours()}: draw detected contours on image.
\end{itemize}

\subsection*{8) NumPy image shape}
\begin{lstlisting}[language=Python]
height, width, channels = img.shape
\end{lstlisting}

\section*{PART II - Detailed Complete Answer}
\subsection*{1) Complete code}
\begin{lstlisting}[language=Python]
import tensorflow as tf
from tensorflow.keras import layers, models

model = models.Sequential([
    layers.Conv2D(16, (3, 3), activation='relu', input_shape=(32, 32, 3)),
    layers.MaxPooling2D((2, 2)),
    layers.Flatten(),
    layers.Dense(3, activation='softmax')
])

model.compile(optimizer='adam', loss='sparse_categorical_crossentropy')
model.summary()
\end{lstlisting}

\subsection*{2) Explanation (line-by-line summary)}
\begin{itemize}[leftmargin=*]
\item Input is RGB 32x32: \texttt{input\_shape=(32,32,3)}.
\item \texttt{Conv2D(16,(3,3),relu)} extracts image features.
\item \texttt{MaxPooling2D((2,2))} downsamples feature maps.
\item \texttt{Flatten()} converts feature maps to a 1D vector.
\item \texttt{Dense(3, softmax)} outputs class probabilities for 3 classes.
\item Compile with Adam + sparse categorical crossentropy as required.
\end{itemize}

\subsection*{3) Common mistakes to avoid}
\begin{enumerate}[leftmargin=*]
\item Wrong input shape (must include channel dimension 3).
\item Wrong output size (must be 3 neurons).
\item Wrong output activation (use softmax for multi-class).
\item Wrong loss function (use sparse\_categorical\_crossentropy).
\end{enumerate}

\section*{PART III - Debugging \& Correction}
\subsection*{Four errors identified}
\begin{enumerate}[leftmargin=*]
\item \texttt{cv2.waitkey(0)} should be \texttt{cv2.waitKey(0)}.
\item \texttt{cv2.release()} is invalid in this image script and should be removed.
\item Missing check after \texttt{cv2.imread("photo.png")}; it may return \texttt{None}.
\item Program should not continue processing if image loading fails.
\end{enumerate}

\subsection*{Corrected code}
\begin{lstlisting}[language=Python]
import cv2

image = cv2.imread("photo.png")

if image is None:
    print("Error: Image not found.")
else:
    gray = cv2.cvtColor(image, cv2.COLOR_BGR2GRAY)
    ret, thresh = cv2.threshold(gray, 128, 255, cv2.THRESH_BINARY)

    cv2.imshow("Original", image)
    cv2.imshow("Threshold", thresh)
    cv2.waitKey(0)

    cv2.imwrite("result.png", thresh)
    cv2.destroyAllWindows()
\end{lstlisting}

\end{document}
